% GENERATED FILE. Edit canonical lessons, then run npm run generate:books.
% canonical-source-hash: fnv1a64:9bfe09f3cfc1d438
% canonical-lessons: ZH-C342-shangyue, ZH-C342-zhongqiu, ZH-W342-qiu-autumn, ZH-C342-qijian, ZH-C342-guihua, ZH-C342-jiaxiang

\chapter{Mid-Autumn}
\label{ch:zh-mid-autumn}
\hlchaptermodality{342}

\begin{chapteropening}
\textbf{By the end of this chapter:} \emph{I can talk about the Mid-Autumn Festival, enjoying the moon, the festival days, osmanthus and home.}

Complete the last lesson of chapter 342: I can talk about the Mid-Autumn Festival, enjoying the moon, the festival days, osmanthus and home.
\end{chapteropening}

\section[shǎngyuè]{shǎngyuè — to enjoy the moon}

\label{lesson:ZH-C342-shangyue}



\begin{quote}
Before the new one: say the Mandarin for to taste food with care, then the Mandarin for trade in a shop.
\end{quote}

\subsection*{You'll want to know: shǎngyuè}
\emph{shǎngyuè} — \textquotedblleft{}to enjoy the moon\textquotedblright{}.

\begin{quote}
\emph{wǎnshang wǒmen yìqǐ zài wàibian shǎngyuè} — In the evening we sat outside together enjoying the moon.
\end{quote}

Its characters come later; for now it is a word for the ear.

\begin{cousinweb}
\emph{shǎngyuè} is built from \emph{shǎng} and \zh{月} 'the moon', both inherited from Middle Chinese, \emph{syangX} and \emph{ngjwot} (Guangyun, in Baxter's transcription), which glosses the first 'to bestow'. The sense 'to enjoy' is not given there.
\end{cousinweb}

\begin{grammarlens}[title={What you've built}]
One more word for telling what happened.
\end{grammarlens}

\subsection*{Your turn}
\begin{itemize}
\raggedright
  \item \emph{Say it:} \emph{shǎngyuè}
  \item \emph{Say it:} \emph{shǎngyuè}, once more
  \item \emph{Say it:} \emph{wǎnshang wǒmen yìqǐ zài wàibian shǎngyuè} — In the evening we sat outside together enjoying the moon.
  \item \emph{Recall:} say the Mandarin for to taste food with care, then the Mandarin for trade in a shop, then say \emph{shǎngyuè} again
\end{itemize}

\begin{tcolorbox}[breakable,colback=teal!4,colframe=teal!35!black,title={Before you move on}]
What does shǎngyuè mean? (\textquotedblleft{}To enjoy the moon\textquotedblright{}.) Say it once more.
\end{tcolorbox}
\section[\texorpdfstring{\zh{中秋} (Zhōngqiū)}{中秋 (Zhōngqiū)}]{\zh{中秋} (Zhōngqiū) — the Mid-Autumn Festival}

\label{lesson:ZH-C342-zhongqiu}



\begin{quote}
Before the new one: say the Mandarin for trade in a shop, then the Mandarin for to enjoy the moon.
\end{quote}

\subsection*{You'll want to know: \zh{中秋}}
\textbf{\zh{中秋}} — \emph{Zhōngqiū} — \textquotedblleft{}the Mid-Autumn Festival\textquotedblright{}.

\begin{quote}
\textbf{\zh{中秋的时候我们都吃月饼}} \emph{Zhōngqiū de shíhou wǒmen dōu chī yuèbǐng} — At the moon festival we all eat mooncakes.
\end{quote}

The new character, in pieces: \zh{秋} — nine strokes: on the left the grain piece of \zh{和}, five strokes; on the right \zh{火}. The next lesson puts it in your hand.

\begin{cousinweb}
\textbf{\zh{中秋}} (\emph{Zhōngqiū}) is built from \zh{中} 'the middle' and \zh{秋} 'autumn', both inherited from Middle Chinese, \emph{trjuwng} and \emph{tshjuw} (Guangyun, in Baxter's transcription), which glosses the first, among other senses, 'half', and says of the second, quoting an older dictionary, that it is when the grain ripens.
\end{cousinweb}

\begin{grammarlens}[title={What you've built}]
One more word for telling what happened.
\end{grammarlens}

\subsection*{Your turn}
\begin{itemize}
\raggedright
  \item \emph{Say it:} \emph{Zhōngqiū}
  \item \emph{Say it:} \emph{Zhōngqiū}, once more
  \item \emph{Say it:} \emph{Zhōngqiū de shíhou wǒmen dōu chī yuèbǐng} — At the moon festival we all eat mooncakes.
  \item \emph{Recall:} say the Mandarin for trade in a shop, then the Mandarin for to enjoy the moon, then say \emph{Zhōngqiū} again
\end{itemize}

\begin{tcolorbox}[breakable,colback=teal!4,colframe=teal!35!black,title={Before you move on}]
What does \zh{中秋} mean? (\textquotedblleft{}The Mid-Autumn Festival\textquotedblright{}.) Say it once more.
\end{tcolorbox}
\section[\texorpdfstring{\zh{秋} (qiū)}{秋 (qiū)}]{\zh{秋} — the grain piece, then \zh{火}}

\label{lesson:ZH-W342-qiu-autumn}



\begin{quote}
The word first, then the character inside it.

\begin{itemize}
\raggedright
  \item \emph{Recall:} say \emph{Zhōngqiū} — \textbf{R1}, one lesson back, before it decays
\end{itemize}

\begin{quote}
\zh{和} \zh{火}
\end{quote}

\emph{Read it:} each one aloud, with nothing said first — \emph{hé}, \emph{huǒ}
\end{quote}

\subsection*{Script — the grain piece, then \zh{火}}
\hlblockfigure{\detokenize{figures/ZH-W342-qiu-autumn-filmstrip.pdf}}{How \zh{秋} is written, stroke by stroke}

\begin{quote}
\zh{秋}
\end{quote}

\textbf{Two pieces, left before right.}

\textbf{The left piece, the grain piece, five strokes,} as on the left of \zh{和}: a short stroke falling to the left at the top, a bar, a vertical down through it, a stroke falling to the left from the crossing, and a dot to the right.

\textbf{The right piece, \zh{火}, four strokes,} as you wrote it: a short stroke on the left, a short stroke on the right falling to the left, then a long stroke falling to the left and a long stroke falling to the right.

5 + 4 = nine.

\subsection*{Your turn}
\begin{itemize}
\raggedright
  \item \emph{Write it:} \zh{秋}, all nine, with the model in front of you
  \item \emph{Cover:} write it again from memory, and count the pieces, not the strokes
\end{itemize}

\begin{tcolorbox}[breakable,colback=teal!4,colframe=teal!35!black,title={Before you move on}]
How many strokes in \zh{秋}? (\textbf{Nine}.) Where have you met the left piece? (\textbf{On the left of \zh{和}}.) What is the right piece? (\textbf{\zh{火}}.)

Source: \href{https://raw.githubusercontent.com/chanind/hanzi-writer-data/68d10a4b21150cae5e1ebbd223eed289cf32d90c/data/\%E7\%A7\%8B.json}{Hanzi Writer Data \zh{秋}, pinned stroke paths and medians}.
\end{tcolorbox}
\section[\texorpdfstring{\zh{期间} (qījiān)}{期间 (qījiān)}]{\zh{期间} (qījiān) — period (during which); during}

\label{lesson:ZH-C342-qijian}



\begin{quote}
Before the new one: say the Mandarin for to enjoy the moon, then the Mandarin for the Mid-Autumn Festival.
\end{quote}

\subsection*{You'll want to know: \zh{期间}}
\textbf{\zh{期间}} — \emph{qījiān} — \textquotedblleft{}period (during which); during\textquotedblright{}.

\begin{quote}
\textbf{\zh{中秋期间他给家人买了月饼}} \emph{Zhōngqiū qījiān tā gěi jiārén mǎi le yuèbǐng} — Over the moon festival he bought mooncakes for his family.
\end{quote}

\begin{cousinweb}
\textbf{\zh{期间}} (\emph{qījiān}) is built from \zh{期} 'a set time' and \zh{间} 'between', both inherited from Middle Chinese, \emph{gi} and \emph{kean} (Guangyun, in Baxter's transcription); the dictionary lists the second under an older form of the character and glosses it 'a gap'. Middle Chinese \emph{gi} would regularly give qí; the level tone of qī is irregular.
\end{cousinweb}

\begin{grammarlens}[title={What you've built}]
One more word for telling what happened.
\end{grammarlens}

\subsection*{Your turn}
\begin{itemize}
\raggedright
  \item \emph{Say it:} \emph{qījiān}
  \item \emph{Say it:} \emph{qījiān}, once more
  \item \emph{Say it:} \emph{Zhōngqiū qījiān tā gěi jiārén mǎi le yuèbǐng} — Over the moon festival he bought mooncakes for his family.
  \item \emph{Recall:} say the Mandarin for to enjoy the moon, then the Mandarin for the Mid-Autumn Festival, then say \emph{qījiān} again
\end{itemize}

\begin{tcolorbox}[breakable,colback=teal!4,colframe=teal!35!black,title={Before you move on}]
What does \zh{期间} mean? (\textquotedblleft{}Period (during which); during\textquotedblright{}.) Say it once more.
\end{tcolorbox}
\section[guìhuā]{guìhuā — osmanthus flowers}

\label{lesson:ZH-C342-guihua}



\begin{quote}
Before the new one: say the Mandarin for the Mid-Autumn Festival, then the Mandarin for the time during an event.
\end{quote}

\subsection*{You'll want to know: guìhuā}
\emph{guìhuā} — \textquotedblleft{}osmanthus flowers\textquotedblright{}.

\begin{quote}
\emph{qiūtiān de shíhou guìhuā kāi le} — In the autumn the osmanthus came into flower.
\end{quote}

Its characters come later; for now it is a word for the ear.

\begin{cousinweb}
\emph{guìhuā} is built from \emph{guì} and \zh{花} 'a flower', both inherited from Middle Chinese, \emph{kwejH} and \emph{xwae} (Guangyun, in Baxter's transcription), which describes the first as a tree that grows in groves and stays green in winter and summer.
\end{cousinweb}

\begin{grammarlens}[title={What you've built}]
One more word for telling what happened.
\end{grammarlens}

\subsection*{Your turn}
\begin{itemize}
\raggedright
  \item \emph{Say it:} \emph{guìhuā}
  \item \emph{Say it:} \emph{guìhuā}, once more
  \item \emph{Say it:} \emph{qiūtiān de shíhou guìhuā kāi le} — In the autumn the osmanthus came into flower.
  \item \emph{Recall:} say the Mandarin for the Mid-Autumn Festival, then the Mandarin for the time during an event, then say \emph{guìhuā} again
\end{itemize}

\begin{tcolorbox}[breakable,colback=teal!4,colframe=teal!35!black,title={Before you move on}]
What does guìhuā mean? (\textquotedblleft{}Osmanthus flowers\textquotedblright{}.) Say it once more.
\end{tcolorbox}
\section[jiāxiāng]{jiāxiāng — hometown; native place}

\label{lesson:ZH-C342-jiaxiang}



\begin{quote}
Before the new one: say the Mandarin for the time during an event, then the Mandarin for sweet osmanthus flowers.
\end{quote}

\subsection*{You'll want to know: jiāxiāng}
\emph{jiāxiāng} — \textquotedblleft{}hometown; native place\textquotedblright{}.

\begin{quote}
\emph{Zhōngqiū de wǎnshang wǒ hěn xiǎng jiāxiāng} — On the night of the moon festival I miss home very much.
\end{quote}

Its characters come later; for now it is a word for the ear.

\begin{cousinweb}
\emph{jiāxiāng} is built from \zh{家} 'a home' and \emph{xiāng} 'a village', both inherited from Middle Chinese, \emph{kae} and \emph{xjang} (Guangyun, in Baxter's transcription), which glosses the second 'a village community'.
\end{cousinweb}

\begin{grammarlens}[title={What you've built}]
One more word for telling what happened.
\end{grammarlens}

\subsection*{Your turn}
\begin{itemize}
\raggedright
  \item \emph{Say it:} \emph{jiāxiāng}
  \item \emph{Say it:} \emph{jiāxiāng}, once more
  \item \emph{Say it:} \emph{Zhōngqiū de wǎnshang wǒ hěn xiǎng jiāxiāng} — On the night of the moon festival I miss home very much.
  \item \emph{Recall:} say the Mandarin for the time during an event, then the Mandarin for sweet osmanthus flowers, then say \emph{jiāxiāng} again
\end{itemize}

\begin{tcolorbox}[breakable,colback=teal!4,colframe=teal!35!black,title={Before you move on}]
What does jiāxiāng mean? (\textquotedblleft{}Hometown; native place\textquotedblright{}.) Say it once more.
\end{tcolorbox}
